\section{The holographic arithmetic-geometric coupling constant}
\label{lib01:acoplamiento}

The dodecaphasic register links two realizations of the same state: the positional reading that enters into arithmetic closure and the reading that preserves incidence, orientation, and direction. The following development determines the operations made possible by the register already produced, the information retained by its memories, and the errors under which its relations can be recovered.
Its reference function admits a precise formulation: the phase orbit of the register is a complete basis of its carrier, and isometric memory preserves exactly the Gram matrix of that basis.

The starting point remains the composition
\[
 \APP\longrightarrow\mathrm{TRIT}\longrightarrow\TPK
 \longrightarrow\text{enriched state}
 \longrightarrow\text{joint discrete structure of the continuum}.
\]
APP supplies the sheets and their remainders and quotients; TRIT preserves regime and orientation; TPK selects, transports, and updates while preserving carry, route, boundary, survival, and memory. The five constructions of the continuum belong to this joint production, developed in the preceding body. The work here takes place after it, on its dodecaphasic integer representation and its typed readers. Neither the subsequent linear carrier nor a decoder's candidate set replaces the domain of generated histories.

\subsection{Object, chart, and production of the register}
\label{lib01:registro}

Definition~\ref{lib01:definicion} fixes the generated object and its readers. We now develop its reversible encoding, its capacity for operator identification, and the stability of its incidence realizations.

The designation expresses a constitutive function of the object, not an automatic identification with a physical interaction coefficient. Canonicity refers to the declared profile, origin, and orientation. A chart transformation and an update that incorporates new events are not the same operation. Nor is it asserted that the same numerical vector must be a necessary invariant of every system.

The calendar of one hundred and eight positions consists of twelve nonadic windows. The extension
\[
 0\longrightarrow C_{12}\longrightarrow C_{108}
 \longrightarrow C_9\longrightarrow0
\]
preserves phase carry; its non-splitting and the distinction between phase return and state return are proved in \cref{prop:k-extension}. In each window, the register of a history retains the subroute, its orientation, the carry, and the incidence register. The causal order of extraction is
\begin{equation}
 x_{\rm term}\longmapsto\mathcal R_{12}(x_{\rm term})
 \longmapsto(b^{90},b^{120},Q_{\rm TPK})
 \longmapsto U_{\rm sgn}\longmapsto K.
 \label{lib01:cadena-registro}
\end{equation}
The regional coordinates and the signed register are two readings of the same terminal state; it is not postulated that the visible ternary matrix alone determines the signed coordinate.

To specify the final composition explicitly, the channels produced are
\begin{align*}
 b^{90}&=(-495,-116,-684,108,601,-90,
                -173,-88,313,560,-397,461),\\
 b^{120}&=(-425,-281,-481,671,-255,30,
                 475,-543,680,251,6,-128),\qquad Q_{\rm TPK}=6263.
\end{align*}
Let \(B_r=\sum_{j=0}^3b^{120}_{r+3j}\) and \(d_{r,j}=b^{90}_{r+3j}\), for \(r=1,2,3\). The harmonic reader gives
\[
 s_1=\frac{Q_{\rm TPK}+2B_1+B_2}{3},\qquad
 s_2=s_1-B_1,\qquad s_3=s_2-B_2,
\]
\[
 U^{(r)}=(s_r,d_{r,1}-d_{r,3},d_{r,0}+d_{r,2},d_{r,0}-d_{r,2}).
\]
In this compatible output, \((B_1,B_2,B_3)=(972,-1073,101)\), \((s_1,s_2,s_3)=(2378,1406,2479)\), and
\begin{align*}
 U^{(1)}&=(2378,-452,-668,-322),\\
 U^{(2)}&=(1406,998,-204,-28),\\
 U^{(3)}&=(2479,-551,-371,-997).
\end{align*}
Division by three belongs to the reader's three orbits, and the subsequent division by four to the Hadamard inverse; these are not coefficients fitted to a physical output. On the orbits \((1,4,7,10)\), \((2,5,8,11)\), \((3,6,9,12)\), we apply
\[
 H_4=\begin{pmatrix}
 1&1&1&1\\1&1&-1&-1\\1&-1&1&-1\\1&-1&-1&1
 \end{pmatrix},\qquad H_4^2=4I.
\]
The products \(H_4U^{(r)}\) are, respectively,
\[
 (936,2916,2484,3176),\quad(2172,2636,232,584),
 \quad(560,3296,3656,2404).
\]
Their divisibility by four verifies membership in the Hadamard sublattice. After dividing and restoring the natural order, we obtain
\begin{equation}
 K=(234,543,140,729,659,824,621,58,914,794,146,601),
 \qquad\mathbf1^*K=6263.
 \label{lib01:K}
\end{equation}
The inversion and its congruences, including the index \(4096\) of the signed image, are proved in \cref{lem:k-hadamard,prop:k-sello}. Reconstruction is unique because \(H_4^{-1}=H_4/4\), not because a register close to a known value of alpha was sought.

To avoid a collision of symbols, we shall write \(\Delta_j=I-S^j\) for the antecedent integer extractors, with \((Sx)_i=x_{i+1}\) and indices modulo twelve. Then
\[
 \Delta_3K=b^{90},\qquad \Delta_4K=b^{120},\qquad Q(K)=6263.
\]
The extractor \((\Delta_3,\Delta_4,Q)\) is injective, with invariant factors \(1,\ldots,1,12\), according to \cref{prop:k-transversal}. Its integral-image conditions and accumulation law remain those of the antecedent register. The normalized memory complements to be denoted by \(D_j\) are different operators, derived from those shifts and with a declared metric codomain.

\subsection{Periodic encoding, closure, and exact inversion}
\label{lib01:kappa}

Fix \(B=1000\), twelve positions, and
\[
 N(k)=\sum_{j=1}^{12}k_jB^{12-j},\qquad d_B=B^{12}-1,
 \qquad \kappa(k)=\frac{N(k)}{d_B}.
\]
In the canonical chart,
\begin{equation}
 \begin{aligned}
 \kappa_{\rm ag}
 &=\frac{234543140729659824621058914794146601}{10^{36}-1},\\
 \alpha_A&=A_{reg}-\kappa_{\rm ag},\qquad
 A_{reg}=\pi_{\rm HMT}+e_{\rm HMT}-\varphi_{\rm HMT}-4.
 \end{aligned}
 \label{lib01:clausura}
\end{equation}
The three regional coordinates are expressed before acting in this reader; their common origin with the dodecaphasic closure does not require conflating the evaluation order of the correlated coordinates. The equality is a subsequent reader of the generated state. Rearranging it once its outputs are known is a legitimate inverse operation, distinct from using a target value to produce the register.

\begin{proposition}[Reversibility and radius of the encoding]
\label{lib01:inversion-kappa}
On the domain \(\{0,\ldots,999\}^{12}\), the fraction \(\kappa(k)\), with declared base and length, determines \(k\) exactly. The same holds for the exact pair \((A_{reg},\alpha_A)\) when it satisfies \eqref{lib01:clausura}. If
\[
 |\widetilde\kappa-\kappa(k)|<\frac1{2d_B},
\]
rounding \(d_B\widetilde\kappa\) to the nearest integer recovers \(N(k)\) without error. If the two terms of the closure are approximated separately, it suffices that the sum of their absolute errors be less than \(1/(2d_B)\).
\end{proposition}
\begin{proof}
Multiplication by \(d_B\) recovers the integer \(N(k)\). Twelve Euclidean divisions by \(B\), retaining the remainders and the number of positions, recover its blocks, including leading zero blocks. The error hypothesis places the true integer at a distance strictly less than half a unit. For the closure, the triangle inequality is applied to the error in the difference. Two consecutive integers show that the radius cannot be increased uniformly over the entire code domain: their midpoint is ambiguous.
\end{proof}

If a reduced fraction \(p/q\) is received, one computes \(N=d_Bp/q\) and checks integrality and \(0\leq N\leq d_B\). The reduced denominator does not replace the base, length, and chart metadata. The code contains \(10^{36}\) words and requires \(\lceil\log_2(10^{36})\rceil=120\) bits in a fixed-length binary representation. Twelve ten-bit blocks also occupy 120 bits: no additional compression, authentication, or cryptographic integrity has been proved. An exact fraction or the integer \(N\) avoids losing the precision needed for this inversion.

The periodic reading is not to be confused with the truncation \(N/B^{12}\). If \(\kappa_j=\kappa(S^jK)\), positional division gives
\begin{equation}
 \kappa_{j+1}=1000\kappa_j-K_{j+1},\qquad
 K_{j+1}=1000\kappa_j-\kappa_{j+1},\qquad
 \sum_{j=0}^{11}\kappa_j=\frac{6263}{999}.
 \label{lib01:orbita-kappa}
\end{equation}
Thus the scalar orbit preserves all phase differences and makes it possible to read \(B_K=\{j+1:1000\kappa_j-\kappa_{j+1}\geq729\}\). The reversibility of these representations does not imply that a scalar reconstructs the entire enriched history that produced the register.

\subsection{The complete orbit as a reference for its carrier}
\label{lib01:orbita}

In \(V=\mathbb R^{12}\), with its canonical Euclidean inner product, consider
\[
 O_K=(K,SK,\ldots,S^{11}K):\mathbb R^{12}\longrightarrow V.
\]
The shift \(S\) acts on positions of the register. Periodicity alone does not identify it with the full TPK evolution or with a subsequent exceptional symmetry.

\begin{theorem}[Complete cyclic reference and exact Gram matrix]
\label{lib01:marco-ciclico}
The matrix \(O_K\) is invertible. Its Gram matrix has the eigenvalues
\[
\begin{array}{c|c}
 \text{value}&\text{multiplicity}\\\hline
 39225169=6263^2&1\\
 1248025-345420\sqrt3&2\\
 589609&2\\
 1407529&2\\
 1053157&2\\
 1248025+345420\sqrt3&2\\
 697225&1
\end{array}
\]
and satisfies
\begin{equation}
 589609\|c\|^2\leq\|O_Kc\|^2\leq39225169\|c\|^2,
 \quad
 \|O_K^{-1}\|=\frac1{\sqrt{589609}},\quad
 \operatorname{cond}(O_K)=\frac{6263}{\sqrt{589609}}.
 \label{lib01:cotas-orbita}
\end{equation}
The orbit of the centred register has rank eleven.
\end{theorem}
\begin{proof}
The entries of the Gram matrix are the autocorrelations \(\langle K,S^jK\rangle\). In the order \(j=0,\ldots,11\), their values are
\begin{align*}
 (&4251257,2999323,3163385,3287920,3216553,3344743,\\
  &2950064,3344743,3216553,3287920,3163385,2999323).
\end{align*}
Since \(S\) is orthogonal, the Gram matrix is circulant. Substituting the twelve characters \(\exp(2\pi i m/12)\) into that row gives the real table above. All values are positive. The smaller of the two values involving \(\sqrt3\) exceeds \(589609\), because \(658416>345420\sqrt3\), an inequality verified by squaring its two positive sides. The other values in the table are greater than or equal to \(589609\). The uniform mode has value \((\sum K_i)^2=6263^2\); the triangle inequality for the Fourier transform, using \(K_i\geq0\), bounds all the others by that same value. The spectral theorem gives \eqref{lib01:cotas-orbita}. Centring the register removes only the uniform mode; the remaining eleven retain the nonzero values in the table.
\end{proof}

The determinant, with the column order and the orientation of \(S\) fixed here, is
\[
 \det O_K=5483119259214020183850397930008283625.
\]
The absence of a proper period of the vector would not, by itself, prove invertibility: the decisive fact is that none of its Fourier modes vanishes. The condition number is approximately \(8.15643\). If \(K\) is normalized to unit norm, the extreme eigenvalues of the Gram matrix are divided by \(\|K\|^2=4251257\); the amplitude of its components must not be confused with an intrinsic improvement over every reference of equal norm.

\begin{corollary}[Polar separation of form and orientation]
\label{lib01:polar}
Let
\[
 G_K=O_K^*O_K,\qquad P_K=G_K^{1/2},\qquad
 F_K=O_KG_K^{-1/2}.
\]
Then \(P_K\) is positive definite, \(F_K\) is orthogonal, and \(O_K=F_KP_K\). An isometric realization \(J:V\to\mathcal Y\) preserves \(G_K\) and \(P_K\), and transports the isometric part to \(JF_K\).
\end{corollary}
\begin{proof}
Positivity of the Gram matrix allows its square root and inverse to be defined through the spectral theorem. The identity \(F_K^*F_K=G_K^{-1/2}G_KG_K^{-1/2}=I\) proves the assertion. Moreover, \((JO_K)^*(JO_K)=O_K^*O_K\).
\end{proof}

This polar decomposition is defined from the complete Gram matrix: it is not obtained solely from the list of its eigenvalues or from the norm of \(K\). The symbol \(P_K\) in this corollary is not the rank-ten dimensional projector to be constructed below. This separation is a consequence of the cyclic reference, not an identification of all its readers.

\subsection{Operator identification and relative holonomy}
\label{lib01:identificacion}

\begin{theorem}[Twelve general responses or one cyclic response]
\label{lib01:operadores}
For any linear operator \(H:V\to V\), its twelve responses \(Y=(HK,HSK,\ldots,HS^{11}K)\) determine
\[
 H=YO_K^{-1}.
\]
If \(HS=SH\), a single complete vector-valued response \(y=HK\) determines
\begin{equation}
 H=O_yO_K^{-1},\qquad
 H=\sum_{j=0}^{11}h_jS^j,\qquad h=O_K^{-1}y.
 \label{lib01:filtro}
\end{equation}
With errors \(E\) in the response matrix or \(\delta y\) in the vector, respectively,
\begin{align}
 \|\widehat H-H\|&\leq\frac{\|E\|}{\sqrt{589609}},
 \label{lib01:error-matriz}\\
 \|\widehat h-h\|&\leq\frac{\|\delta y\|}{\sqrt{589609}},\qquad
 \|\widehat H-H\|\leq\sqrt{\frac{12}{589609}}\,\|\delta y\|.
 \label{lib01:error-filtro}
\end{align}
\end{theorem}
\begin{proof}
The first equality is \(Y=HO_K\). If \(H\) commutes with \(S\), then \(HS^jK=S^jy\), so that \(Y=O_y\). Commutation with the cycle determines all columns of \(H\) from any one of them; the twelve powers \(I,S,\ldots,S^{11}\) are independent and span that space of operators. Hence \(y=O_Kh\), which gives the second formula. The first two bounds use \(\|O_K^{-1}\|=1/\sqrt{589609}\). The last uses \(\|O_{\delta y}\|\leq\|O_{\delta y}\|_{\rm F}
=\sqrt{12}\|\delta y\|\).
\end{proof}

A response here means a vector of twelve components, not a scalar measurement. For example, if \(H=2I+3S-S^5\), the response \(y=HK\) returns exactly the coefficients \(2,3,-1\) in their positions through \eqref{lib01:filtro}; the inversion procedure needs \(y,K,S\), not the filter coefficients. By contrast, the nonzero operator whose first row is \((543,-234,0,\ldots,0)\) and whose remaining rows are zero annihilates \(K\). A single response does not identify an arbitrary operator.

The reference function admits an exact generalization. On a cycle of \(n\) positions, \(H\mapsto Hg\) identifies all operators commuting with the cycle if and only if \(O_g=(g,Sg,\ldots,S^{n-1}g)\) is invertible. Sufficiency is the preceding proof. If \(O_g\) is singular, a nonzero vector in its kernel gives a polynomial of degree less than \(n\) that annihilates \(g\); the operator is nonzero because the powers of the cycle are independent. The unit impulse also has an orthonormal orbit and serves as a reference. The specific consequence for \(K\) is that the same object already used in closure and incidence satisfies this condition, without having been selected anew to optimize the test. No identification of nonlinear dynamics or arbitrary hidden states follows; an application to a linearization requires that the linearization have been defined and justified.

\begin{corollary}[Recovery of relative holonomy from the complement]
\label{lib01:holonomia}
Let \(H_B\) be known and invertible, and let \(H_C\) be another transport of the same type. The complement
\[
 D_\gamma=\frac{\sqrt8}{9}(H_C-H_B)
\]
determines
\[
 H_B^{-1}H_C=I+\frac9{\sqrt8}H_B^{-1}D_\gamma.
\]
If \(D_\gamma\) commutes with \(S\), in particular if both transports do, \(y=D_\gamma K\) suffices to recover
\begin{equation}
 H_B^{-1}H_C
 =I+\frac9{\sqrt8}H_B^{-1}O_yO_K^{-1}.
 \label{lib01:holonomia-una-respuesta}
\end{equation}
The error in this reconstruction is bounded by
\[
 \frac9{\sqrt8}\|H_B^{-1}\|
 \sqrt{\frac{12}{589609}}\,\|\delta y\|.
\]
\end{corollary}
\begin{proof}
Solve for \(H_C\) in the definition of the complement and apply \cref{lib01:operadores}. The bound follows from submultiplicativity.
\end{proof}

For \(H_B=S^3\) and \(H_C=S^4\), the rational datum \(y/\sqrt8=(S^4-S^3)K/9\) makes it possible to reconstruct the relative holonomy \(S\) without supplying \(H_C\) to the inversion procedure. If \(H_B\) is unitary, the error factor is approximately \(0.01435510\). Without commutation, the twelve general responses are retained, rather than a fictitious symmetry hypothesis. The selection of a physical loop and its representation still belong to its generating dynamics; this protocol does not replace them.

\subsection{Complete memory preserves the reference and the observables}
\label{lib01:memoria-completa}

Let \(C_n\) be unitary operators on a Hilbert space \(H\),
\[
 T_n=\frac{8I+C_n}{9},\qquad
 D_n=\frac{\sqrt8(C_n-I)}9,\qquad
 R_0=I,\quad R_{n+1}=T_nR_n.
\]
This formulation includes nonstationary sequences. It does not require commutation between \(C_n\) at different stages.

\begin{theorem}[Isometric register, terminal component, and transport]
\label{lib01:isometria}
For every \(N\),
\[
 R_N^*R_N+\sum_{n=0}^{N-1}R_n^*D_n^*D_nR_n=I.
\]
The positive strong limit \(A_\infty=\lim_N R_N^*R_N\) exists, and
\begin{equation}
 \mathcal Vx=
 \bigl(A_\infty^{1/2}x,(D_nR_nx)_{n\geq0}\bigr)
 \quad\text{satisfies}\quad\mathcal V^*\mathcal V=I.
 \label{lib01:V-infinita}
\end{equation}
Its image \(\mathcal M\) is closed, and the inverse on it is \(\mathcal V^*\), of norm one. In the finite case, the first component is replaced by \(R_Nx\) and the \(N\) complements are retained. For any bounded operator \(A\),
\[
 A^{\mathcal V}=\mathcal VA\mathcal V^*\big|_{\mathcal M}
\]
preserves products, adjoints, norm, and spectrum on that image, and
\[
 \langle\mathcal Vx,A^{\mathcal V}\mathcal Vy\rangle
 =\langle x,Ay\rangle.
\]
\end{theorem}
\begin{proof}
Expanding the two terms of one stage gives \(T_n^*T_n+D_n^*D_n=I\). Hence \(R_n^*R_n-R_{n+1}^*R_{n+1}=R_n^*D_n^*D_nR_n\), and the telescoping sum proves the finite identity. The sequence of positive operators \(R_N^*R_N\) is decreasing and bounded below; its strong limit exists. Taking inner products with \(x\), the sum of the squared norms of the complements is \(\|x\|^2-\langle x,A_\infty x\rangle\). This proves that the series belongs to the Hilbert direct sum and that \(\mathcal V\) is an isometry. An isometry has closed image and an adjoint inverse on it. Inserting \(\mathcal V^*\mathcal V=I\) into the compositions proves the transport properties. Neither convergence of \(R_Nx\) nor the assertion that \(A_\infty\) is a projector has been assumed.
\end{proof}

Applying this result to the carrier of the register,
\begin{equation}
 (\mathcal VO_K)^*(\mathcal VO_K)=O_K^*O_K.
 \label{lib01:gram-memoria}
\end{equation}
The twelve directions and all their correlations are preserved, with the same bounds as in \cref{lib01:marco-ciclico}, at any depth of complete memory. It is not asserted that the entire continuum has dimension twelve: this twelve-dimensional reference is preserved in its image. The correct operators there are \(S^{\mathcal V}=\mathcal VS\mathcal V^*\) and \(H^{\mathcal V}=\mathcal VH\mathcal V^*\). An ambient shift of storage slots is not presumed to be \(S^{\mathcal V}\).

If \(z=\mathcal Vy+\delta z\), the estimate
\[
 \widehat h=O_K^{-1}\mathcal V^*z
\]
retains the first bound of \eqref{lib01:error-filtro} with \(\|\delta z\|\) in place of \(\|\delta y\|\). Depth does not worsen this constant because the isometry and its terminal component are preserved. A partial selection of channels has a different Gram matrix and requires its own bounds. The exterior powers \(\bigwedge^p\mathcal V\) preserve Gram determinants; finite-rank traces and spectral measures of the transported observables are also preserved. Squared norms are not thereby interpreted automatically as Born probabilities: that reading requires its subsequent realization.

\subsection{Exceptional incidence and selected direction}
\label{lib01:direccion}

Let \(e=\mathbf1\), \(\Pi=I-ee^*/12=P_{11}\). In the icosahedral chart \(A_5/C_5\) already constructed, the orthogonal representation \(\rho\) and the three-dimensional character fix
\begin{equation}
 P_3=\frac3{60}\sum_{g\in A_5}\chi_3(g^{-1})\rho(g),
 \quad P_3^*=P_3=P_3^2,\quad
 \operatorname{rank}P_3=3,\quad P_3e=0.
 \label{lib01:P3}
\end{equation}
This projector is received together with its representation and markings; it is not chosen to fit a direction to the register. Put
\[
 k=\Pi K,\qquad u=P_3k,\qquad
 E_u=\frac{uu^*}{\|u\|^2},\qquad Q_K=\Pi-E_u=P_{10}(K).
\]
The exact evaluations are
\begin{align}
 \|k\|^2&=\frac{11789915}{12},&
 \|u\|^2&=\frac{6638585+2275584\sqrt5}{20},
 \label{lib01:normas-direccion}\\
 \eta_K:=\frac{\|u\|^2}{\|k\|^2}
 &=\frac{3(6638585+2275584\sqrt5)}{5\cdot11789915}
 \simeq0.5967954228259091.
 \label{lib01:eta}
\end{align}
In particular, \(u\ne0\). Since \(u^*k=\|u\|^2\),
\begin{equation}
 E_uk=u,\qquad Q_Kk=k-u,\qquad
 k=u+Q_Kk,\qquad \|k\|^2=\|u\|^2+\|Q_Kk\|^2.
 \label{lib01:particion-direccion}
\end{equation}
The reduction \(12\to11\to10\) removes first the uniform direction and then a selected line. It does not remove all of \(\operatorname{im}P_3\): in \(\operatorname{im}Q_K\), the two directions of that subspace orthogonal to \(u\) remain. These dimensions are those of the algebraic spaces written above; identifying them with physical dimensions requires their realization maps.

\begin{proposition}[Exact transport of direction and descriptor]
\label{lib01:transporte-direccion}
For a complete isometry \(\mathcal V\), let \(\Pi^{\mathcal V}=\mathcal V\Pi\mathcal V^*\), \(P_3^{\mathcal V}=\mathcal VP_3\mathcal V^*\), and \(u^{\mathcal V}=\mathcal Vu\). On \(\mathcal M\),
\[
 Q_K^{\mathcal V}
 =\mathcal VQ_K\mathcal V^*
 =\Pi^{\mathcal V}
 -\frac{u^{\mathcal V}(u^{\mathcal V})^*}{\|u^{\mathcal V}\|^2},
 \quad \operatorname{rank}Q_K^{\mathcal V}=10,
 \quad \eta_K^{\mathcal V}=\eta_K.
\]
Moreover, \(\mathcal V^*u^{\mathcal V}=u\) and \(\mathcal V^*Q_K^{\mathcal V}\mathcal V=Q_K\).
\end{proposition}
\begin{proof}
Substitute the definitions and \(\mathcal V^*\mathcal V=I\). Norms and rank are preserved by the isometry. The identity on the image is \(\mathcal V\mathcal V^*\), not the identity of a larger ambient storage space. Extension by zero outside the image adds no physical directions.
\end{proof}

The complete incidence screen provides another exact realization. If \(\mathcal I\) is the incidence operator already constructed and \(E_{\rm inc}=\mathcal I(e^\perp)\), then
\begin{equation}
 F:V\longrightarrow Y=\mathbb R\oplus E_{\rm inc},\quad
 Fx=\left(\mu(x),\frac{\mathcal I\Pi x}{6}\right),\quad
 \mu(x)=\frac{\mathbf1^*x}{12},\qquad
 \langle(\mu,y),(\nu,z)\rangle_Y=12\mu\nu+\langle y,z\rangle
 \label{lib01:incidencia-F}
\end{equation}
is a surjective isometry, since \(\Pi\mathcal I^*\mathcal I\Pi=36\Pi\). For the raw incidence of the 132 hexads, the complete identity is \(\mathcal I^*\mathcal I=36I+30\mathbf1\mathbf1^*\); centring removes exactly its uniform component. The incidence image has dimension eleven; \(Y\) is not all of \(\mathbb R\oplus\mathbb R^{132}\). Applying \(F\) term by term to memory preserves both the Gram matrix and the preceding operators and their domains. This composition introduces no Monster action on \(K\).

\subsection{Continuous stability and the need to preserve correlations}
\label{lib01:ruido}

\begin{theorem}[Direction error from complete memory]
\label{lib01:error-proyectores}
Let \(z=\mathcal VK+\delta\), \(\|\delta\|\leq\varepsilon\), and reconstruct \(\widehat K=\mathcal V^*z\), \(\widehat k=\Pi\widehat K\), \(\widehat u=P_3\widehat k\). All three errors are less than or equal to \(\varepsilon\). If \(\varepsilon<\|u\|\),
\[
 \|Q_{\widehat K}-Q_K\|\leq\frac\varepsilon{\|u\|}.
\]
If \(\varepsilon<\|k\|\),
\[
 |\eta_{\widehat K}-\eta_K|\leq\frac\varepsilon{\|k\|}.
\]
\end{theorem}
\begin{proof}
The adjoint of the isometry and the projectors have norm one. Noise orthogonal to \(\mathcal M\) disappears upon applying \(\mathcal V^*\). For two nonzero lines, the norm of the difference of their projectors is the sine of the angle between them. The distance from \(u\) to the line of \(\widehat u\) does not exceed \(\|u-\widehat u\|\), so this sine is less than or equal to \(\varepsilon/\|u\|\). The hypothesis prevents \(\widehat u=0\).

For the second assertion, normalize \(k\) and \(\widehat k\). The difference of their rank-one projectors has eigenvalues \(\pm\sin\theta\) in their plane. When paired with \(0\leq P_3\leq I\), its absolute value does not exceed \(\sin\theta\): in an eigenvector basis one obtains \(\sin\theta(a-b)\), with \(a,b\in[0,1]\). The same distance argument gives \(\sin\theta\leq\varepsilon/\|k\|\). If the lines coincide, the difference of the ratios is zero.
\end{proof}

Here \(\|u\|\simeq765.733\), so the factor for the projector is approximately \(0.00130594\). In the reader with only two memories, the inverse has norm \(9/4\) and the corresponding factor is \(9/(4\|u\|)\simeq0.00293836\). These are different observed data; the difference is not a cost-free improvement of the same channel. The integer correction in \cref{lib02:radios-memoria} will allow the register to be recovered first without error within its radius, instead of propagating a continuous error to the projector.

\begin{proposition}[Cross blocks are part of the observable]
\label{lib01:correlaciones}
Write \(\mathcal Vx=(W_ax)_a\), including the terminal component. Then the \((a,b)\) block of \(A^{\mathcal V}\) is
\[
 (A^{\mathcal V})_{ab}=W_aAW_b^*.
\]
Its quadratic form is computed as a limit of finite truncations
\[
 \langle y,A^{\mathcal V}y\rangle
 =\lim_F\sum_{a,b\in F}\langle y_a,W_aAW_b^*y_b\rangle.
\]
General conservation is not obtained by discarding all terms \(a\ne b\).
\end{proposition}
\begin{proof}
The block formula follows from multiplying \(\mathcal VA\mathcal V^*\). The truncation projectors converge strongly to the identity and \(A^{\mathcal V}\) is bounded, proving the limit without asserting absolute convergence of an arbitrary double series. For a counterexample, take \(H=\mathbb R\), \(C=-I\), \(T=7/9\), \(D=-2\sqrt8/9\), \(\mathcal Vx=(Tx,Dx)\), and \(A=I\). The diagonal part, evaluated at \(\mathcal Vx\), equals
\[
 (T^4+D^4)x^2=\frac{3425}{6561}x^2.
\]
The cross terms contribute \(3136x^2/6561\), and the sum is \(x^2\), as required by the isometry.
\end{proof}

\subsection{Discrete correction before incidence selection}
\label{lib01:correccion-incidencia}

The canonical two-memory readers are
\[
 T_j=\frac{8I+S^j}{9},\qquad
 D_j=\frac{\sqrt8(S^j-I)}9,\qquad
 Mx=(D_3x,D_4T_3x).
\]
The second memory acts after the first compression; the first complement is not reinjected in this protocol. The general theorem in \cref{lib02:radios-memoria} will prove, without using the register as a target, the strict radii
\[
 r_0=\frac{2\sqrt{146}}{81},\qquad
 r_q=\frac{4\sqrt{73}}{81}.
\]
The first radius recovers \(\Pi K\) from \(MK\), modulo uniform integer translation. The second recovers \(K\) in full when \(q=\mathbf1^*K\) is also preserved. If a real-valued charge is observed with error less than \(1/2\), rounding first recovers its integer charge channel. None of these procedures requires the terminal component \(T_4T_3K\).

The high-block selection in the marked chart is
\begin{equation}
 B_K=\{i:K_i\geq729\}=\{4,6,9,10\}.
 \label{lib01:umbral}
\end{equation}
The fourth component is exactly at the threshold. This circumstance separates continuous stability from discrete stability.

\begin{proposition}[Zero continuous margin and integer restoration]
\label{lib01:umbral-estable}
The direct selector in \eqref{lib01:umbral}, applied to the real-valued reconstruction of the memories, has no positive Euclidean stability radius at \(MK\), even with exact charge. By contrast, prior integer decoding with error less than \(r_q\) restores \(B_K\) exactly without changing the threshold.
\end{proposition}
\begin{proof}
Let \(L\) be the least-squares inverse, satisfying \(LM=\Pi\). For \(t>0\), put
\[
 e_t=-tM\Pi e_4,
 \qquad \widetilde K_t=\frac q{12}\mathbf1+L(MK+e_t)
 =K-t\Pi e_4.
\]
We have \(\|e_t\|\to0\), but the fourth component is \(729-11t/12<729\). For small \(t\), only this membership changes. The second assertion follows by first recovering \(K\) exactly through \cref{lib02:radios-memoria} and then applying the same selector.
\end{proof}

If the remaining marked data \(P_\pi,H_e,H_\varphi,H_{\rm APP}\) are kept exact or transported according to their law, uniqueness of the exceptional construction returns the same flag \(B_K\subset H^-_\alpha\) and the same marked origin. Correcting the register is not equivalent to correcting independent errors in these markings. The witness
\[
 K'=K-e_4+e_6,
 \qquad Q(K')=6263,\qquad B_{K'}=\{6,9,10\}
\]
also belongs to \(\{0,\ldots,999\}^{12}\) and satisfies \(\|MK'-MK\|^2=4672/6561\). The midpoint of the two memories is at distance \(r_q\) from each. This proves that the strict radius is uniformly optimal on the fixed-charge algebraic domain, not that \(K'\) has a constructed HMT terminal history.

Without charge, \(M(K-\mathbf1)=MK\) and the high-block selections differ. The two memories can recover \(Q_K\) and \(\eta_K\) exactly, since these ignore the mean, without recovering that absolute selector. The condition \(P_3\Pi K\ne0\) remains necessary for the projector; integer correction does not remove this defining condition. A concrete finite test uses \(z=MK/\sqrt8\) and adds \(1/10\) to its first component. The error in the norm of \(M\) is \(\sqrt8/10<r_q\). The floor-and-ceiling algorithm in \cref{lib02:decodificador} leaves 924 candidates of charge 6263 and recovers \(K\) and \(B_K\) exactly.

Applying the isometry \eqref{lib01:incidencia-F} gives \(M_{\rm inc}=(F\oplus F)M\). All distances between words and all noise norms are preserved, so the same radii and decoders apply in \(Y\oplus Y\). This is the quantitative connection between memory and incidence; it does not depend on attributing an additional exceptional symmetry to the shift \(S\).

\subsection{Covariance, refinement, and variational charge}
\label{lib01:covarianza}

If \(G\) is orthogonal and \(K'=GK\), \(\Pi'=G\Pi G^*\), and \(P_3'=GP_3G^*\) are transported jointly, then
\[
 u'=Gu,\qquad Q_{K'}'=GQ_KG^*,\qquad \eta_{K'}'=\eta_K.
\]
In a fixed chart, \(K\mapsto aK+b\mathbf1\), \(a\ne0\), leaves \(Q_K\) and \(\eta_K\) invariant; it does not leave the entire object invariant. For example, \(\kappa(K+\mathbf1)-\kappa(K)=1/999\), and the reading \(\alpha_A\) would change by \(-1/999\) if the regions were held fixed. This is an algebraic counterexample to identifying the descriptor with the entire register, not an assertion that this update is admissible as an HMT history.

Nor can one move only the register and arbitrarily keep the projector fixed. The finite formula \eqref{lib01:P3} gives exactly
\[
 \|P_3S^3K\|^2=\frac{369035}{4}-\frac{281861\sqrt5}{10},
 \qquad
 \|P_3S^4K\|^2=\frac{1327717}{4}+\frac{694618\sqrt5}{5}.
\]
These differ from \eqref{lib01:normas-direccion}; in particular, \([P_3,S^3]\) and \([P_3,S^4]\) do not vanish. Transporting \(P_3\) by the same conjugation does preserve \(\eta_K\).

Decoding is also covariant. For a permutation \(R\),
\[
 S'=RSR^{-1},\qquad M'=(R\oplus R)MR^{-1},\qquad
 M'R=(R\oplus R)M.
\]
The lattice, charge, and distances are preserved. In the uniqueness regime, the decoder output transforms by \(R\); outside it, the set of minimizers is transported without arbitrarily breaking a tie. A general orthogonal frame also requires transporting \(\mathbb Z^{12}\) to \(R\mathbb Z^{12}\) and the uniform direction to \(R\mathbf1\). Rounding in the old lattice after changing the chart does not represent the same problem.

\begin{proposition}[Naturality under refinement with local frames]
\label{lib01:refinamiento}
Let \(z\) be a cylinder with successors \(z'\), compatible positive weights \(\sum_{\tau z'=z}\mu(z')=\mu(z)\), and orthogonal frames \(G_z\). The refinement
\[
 \mathcal R_n(e_z\otimes x)=
 \sum_{\tau z'=z}\sqrt{\frac{\mu(z')}{\mu(z)}}
 e_{z'}\otimes G_{z'}G_z^*x
\]
is isometric. For the projector field \(Q_z=G_zQ_KG_z^*\) and its block operator \(Q_n\),
\[
 Q_{n+1}\mathcal R_n=\mathcal R_nQ_n,
 \qquad \mathcal R_n^*Q_{n+1}\mathcal R_n=Q_n.
\]
The identities compose between levels and pass to the inductive limit.
\end{proposition}
\begin{proof}
The successors have orthogonal supports and the weights sum to one; the frame changes are orthogonal. In each successor, \(Q_{z'}G_{z'}G_z^*=G_{z'}G_z^*Q_z\), proving the first square. Multiplication by the adjoint of the refinement gives the second. Compatibility under composition defines the operator on the limit of the isometric inclusions. The same argument applies to \(E_u\).
\end{proof}

The statement preserves the representation of a prefix under compatible refinements. New events may produce another register; temporal constancy of \(K\) does not follow. The ranks of fields on tensor fibres are those of these fibres, not a count of physical dimensions.

\begin{proposition}[A transport charge with an explicit action]
\label{lib01:carga-noether}
Let \(U_n:E_n\to E_{n+1}\) be invertible transports. The auxiliary discrete action with cotangent variables is
\[
 \mathcal A(q,p)=\sum_n p_{n+1}^*(q_{n+1}-U_nq_n).
\]
Finitely supported variations give \(q_{n+1}=U_nq_n\) and \(p_n=U_n^*p_{n+1}\). If \(G_0=I\), \(G_{n+1}=U_nG_n\), \(A_0=Q_K\), and \(A_n=G_nA_0G_n^{-1}\), the quantity
\[
 J_n=p_n^*A_nq_n
\]
is independent of \(n\). The action is invariant under the continuous group \(q_n\mapsto\exp(tA_n)q_n\), \(p_n\mapsto\exp(-tA_n^*)p_n\).
\end{proposition}
\begin{proof}
The indicated equations follow by varying \(p_{n+1}\) and \(q_n\). The definition of \(A_n\) implies \(A_{n+1}U_n=U_nA_n\). Hence
\[
 p_{n+1}^*A_{n+1}q_{n+1}
 =p_{n+1}^*U_nA_nq_n=p_n^*A_nq_n.
\]
The same intertwining holds for the exponentials and makes every term of the action invariant. It is not required that \(A_n\) be antisymmetric: the covector transformation is contragredient.
\end{proof}

For orthogonal transports one may choose \(p_n=q_n\); then \(J_n=\|A_nq_n\|^2\), and the complement provides the remaining charge. In a symplectic realization without a positive form, the pairing written above is preserved and is not called positive energy. If the fibres are archived through isometries \(\mathcal V_n\), the definitions
\[
 \widehat U_n=\mathcal V_{n+1}U_n\mathcal V_n^*,\qquad
 \widehat A_n=\mathcal V_nA_n\mathcal V_n^*,\qquad
 \widehat q_n=\mathcal V_nq_n,
\]
with \(\widehat p_n(y)=p_n(\mathcal V_n^*y)\) preserve the intertwining and the same charge on the images. This composition provides a concrete continuous action and symmetry for a conservation law; it attributes no automatic Noether currents to a finite group, nor does it identify the auxiliary action with every physical action.

\subsection{Encoding, operator identification, and incidence stability}
\label{lib01:alcance}

The contribution assembled here is an operational chain: the generated register has a reversible encoding; its orbit is a complete reference; terminal component and complements preserve this reference and the transported observables; the integer separation of two memories allows incidence and direction to be restored even when the real-valued selector has zero margin. The proofs concerning linearity, projection, and lattices are used as the language of subsequent realization, without changing the production of the object.

Each hypothesis has a witness that prevents its inadvertent extension. An isolated response does not identify an operator that fails to commute with the cycle; only complete memory automatically preserves the original Gram matrix; cross terms cannot be omitted; absolute selection requires charge; the projector requires a nonzero direction; noise bounds depend on normalization; a chart transformation requires transporting the form, lattice, and readers. A perturbation of a reading does not become an admissible history by decree. These delimitations preserve the functions of \(K\), without reducing it to twelve disconnected numbers or promoting a canonical property to universal invariance.
